\cotefile{79}
\scanpage{129}\mspage{65}

\diagram{\begin{tikzcd}[ampersand replacement=\&,column sep=large]
(62)\quad 1 \arrow[r] \& G_0 \arrow[r] \& G \arrow[r,"\delta"] \arrow[d] \& Z' \\
\& \& \mathrm{Cent}(G)=Z' \arrow[ur,swap,"\lambda\mapsto\lambda^2"] \&
\end{tikzcd}}

On voit donc que $\delta(G)$ contient \deleted{$Z'^{2}$}
\[
Z'^{2}=\{\lambda^2\mid\lambda\in Z'\},
\]
mais cela ne signifie pas pour autant que $\delta$ soit surjectif, i.e. que $G$ soit transitif sur l'ens des structures $\Pi$. On voit de suite que l'orbite d'un $i_0\in\unsure{\Pi_S|\Pi_s|\pi}$ sous $G$, correspondant à la structure $A_{i_0}=\unsure{A|A_i|A'|A_0}$ sur $S$, est formée de l'ens des structures $A_\lambda=A_{\lambda\cdot i_0}$ sur $S$ ``équivalentes'' à $A_{i_0}$ (au sens de la p. 61) qui lui sont de plus \uline{isomorphes}. L'axiome (To 5) -- qu'on n'avait pas utilisé jusqu'à maintenant) signifie justement qu'on trouve ainsi tout \deleted{$\unsure{Z}$} $\Pi$, i.e. que (\unsure{63|62|61}) est une suite exacte courte
\[
1\longrightarrow G_0\longrightarrow G\xrightarrow{\delta}Z'\longrightarrow 1\,.\tag*{(63)}
\]

% archive sheet number: 128
